% CLAUDE REPORT::WORSTCASE-BRIDGE-0001::2026-11-10T14:51-08:00
% Scheduled Claude routine run on the shared worst-case derivation objective.
% Additive. Changes no task schedules, agent count, role ownership, scientific
% allocation, or root main.tex ownership region. P versus NP remains OPEN.
\documentclass[11pt]{article}
\usepackage{amsmath,amssymb,amsthm}
\usepackage[margin=1in]{geometry}
\usepackage{hyperref}
\newtheorem{theorem}{Theorem}
\newtheorem{proposition}[theorem]{Proposition}
\newtheorem{corollary}[theorem]{Corollary}
\theoremstyle{definition}
\newtheorem{definition}[theorem]{Definition}
\theoremstyle{remark}
\newtheorem{remark}[theorem]{Remark}
\newcommand{\WR}[1]{\overline{W}_{#1}}
\title{Worst-case minimum derivation cost and the bridge $\mathrm{SAT}\in\mathrm{P}\Rightarrow\WR{R}(N)\le\mathrm{poly}(N)$\\
\large CLAUDE WORSTCASE-BRIDGE-0001}
\date{2026-11-10 14:51 PST}
\begin{document}
\maketitle

\begin{abstract}
The shared objective asks for a superpolynomial lower bound on
\[
\WR{R}(N)=\max_{F\in\mathrm{UNSAT},\ |\langle F\rangle|\le N}\ \min_{\pi:F\vdash_R\bot}\operatorname{cost}_R(\pi)
\]
for a fixed sound and complete calculus $R$, plus the ``weakest sufficient bridge''
$\mathrm{SAT}\in\mathrm{P}\Rightarrow\WR{R}(N)\le\mathrm{poly}(N)$.
We classify this bridge exactly. (i) $\WR{R}\le\mathrm{poly}$ is precisely the statement that $R$ is a
polynomially bounded proof system in the sense of Cook and Reckhow. (ii) $\mathrm{SAT}\in\mathrm{P}$ always makes
\emph{some} proof system polynomially bounded, so the bridge for a fixed $R$ is an optimality claim about $R$.
(iii) For every $R$ for which a superpolynomial $\WR{R}$ is already a theorem (resolution, tree-like resolution, Res$(k)$,
polynomial calculus, cutting planes, bounded-depth Frege) the bridge is \emph{equivalent} to $\mathrm{P}\ne\mathrm{NP}$; it is not weaker.
(iv) For Extended Frege the bridge is a known theorem for polynomial-time algorithms whose correctness is provable in Cook's theory PV,
but the matching lower bound on $\WR{\mathrm{EF}}$ is a major open problem. We call (iii)--(iv) the \emph{bridge dilemma}.
Charging nothing for the choice of the optimal derivation is correct and costs nothing in rigor, but it makes $\WR{R}$ a nondeterministic
(proof-complexity) quantity; the entire difficulty sits in the bridge. \textbf{P versus NP remains OPEN.}
\end{abstract}

\section*{Status summary}
\begin{tabular}{lll}
Item & Statement & Status\\\hline
Prop.~\ref{p:pbounded} & $\WR{R}\le\mathrm{poly}\iff R$ p-bounded & PROVED/EXACT (definitional)\\
Prop.~\ref{p:RA} & $\mathrm{SAT}\in\mathrm{P}\Rightarrow\exists R_A$ with $\WR{R_A}\le\mathrm{poly}$ & PROVED/EXACT\\
Thm.~\ref{t:dilemma}(a) & Bridge$(R)\iff\mathrm{P}\ne\mathrm{NP}$ when $\WR{R}$ superpoly is known & PROVED/EXACT\\
Thm.~\ref{t:dilemma}(b) & Bridge(EF) for PV-verified $A$ & KNOWN THEOREM\\
Open & $\WR{\mathrm{EF}}$ superpolynomial & OPEN\\
Open & Bridge(EF) for all polynomial $A$ & OPEN\\
Prop.~\ref{p:cdcl} & Algorithm-class bridge for CDCL & KNOWN THEOREM\\
Multisign DAG calculus & see \S\ref{s:hyper} & CONDITIONAL on rule set
\end{tabular}

\section{Definitions}
\begin{definition}
A \emph{proof system} for UNSAT (Cook--Reckhow) is a polynomial-time verifier $V(F,\pi)$ such that $F\in\mathrm{UNSAT}$ iff $\exists\pi\,V(F,\pi)=1$.
A cost function $\operatorname{cost}_R$ is \emph{admissible} if $|\pi|\le\operatorname{cost}_R(\pi)\le q(|\pi|)$ for a fixed polynomial $q$
(unit-cost or bounded-weight inference counting with polynomial-size lines satisfies this up to a polynomial).
$R$ is \emph{polynomially bounded} (p-bounded) if every $F\in\mathrm{UNSAT}$ has an $R$-refutation of size $\mathrm{poly}(|\langle F\rangle|)$.
\end{definition}
\begin{definition}
Bridge$(R)$ is the sentence: $\mathrm{SAT}\in\mathrm{P}\Rightarrow\WR{R}(N)\le N^{O(1)}$.
\end{definition}

\section{Exact classification}
\begin{proposition}\label{p:pbounded}
For admissible cost, $\WR{R}(N)\le N^{O(1)}$ iff $R$ is p-bounded.
\end{proposition}
\begin{proof}
The max over $F$ with $|\langle F\rangle|\le N$ of the min over refutations is exactly the worst-case shortest-proof function, up to the polynomial $q$.
\end{proof}

\begin{proposition}\label{p:RA}
If a deterministic algorithm $A$ decides SAT in time $N^k$, then the system $R_A$ with verifier
$V_A(F,\pi)=1\iff\pi=\varepsilon\wedge A(F)=\mathrm{UNSAT}$ is a sound complete proof system with $\WR{R_A}(N)=O(1)$ in proof length and $O(N^k)$ in verification time.
\end{proposition}
\begin{proof}
Immediate. (Cook--Reckhow 1979: NP $=$ coNP iff a p-bounded proof system exists.)
\end{proof}

\begin{remark}
Prop.~\ref{p:RA} shows that the bridge is \emph{trivially true for some} $R$. For a \emph{fixed} $R$, Bridge$(R)$ says that whenever such an $A$ exists,
$R$ also has polynomial refutations, i.e.\ $R$ simulates $R_A$ at least non-effectively. This is an optimality property of $R$
(Krajíček--Pudlák 1989). It is never proved by analyzing $R$ alone: one must show that the run of an \emph{arbitrary} polynomial algorithm
can be turned into a short $R$-derivation.
\end{remark}

\begin{theorem}[Bridge dilemma]\label{t:dilemma}
\begin{enumerate}
\item[(a)] Let $R$ be any proof system for which $\WR{R}(N)=N^{\omega(1)}$ is already a theorem. Then Bridge$(R)\iff\mathrm{P}\ne\mathrm{NP}$.
This applies to resolution (Haken 1985; Urquhart 1987; Chvátal--Szemerédi 1988; Ben-Sasson--Wigderson 2001), tree-like resolution, Res$(k)$ for constant $k$,
polynomial calculus (Razborov 1998; Alekhnovich--Razborov), cutting planes (Pudlák 1997), and bounded-depth Frege (Ajtai 1988; Beame--Impagliazzo--Pitassi; Krajíček--Pudlák--Woods).
\item[(b)] For $R=$ Extended Frege: if PV proves ``$A$ is a polynomial-time SAT decider'', then $\WR{\mathrm{EF}}(N)\le N^{O(1)}$ (Cook 1975 translation; Krajíček--Pudlák 1990).
No superpolynomial lower bound on $\WR{\mathrm{EF}}$ is known.
\end{enumerate}
\end{theorem}
\begin{proof}
(a) ($\Rightarrow$) Bridge$(R)$ and the known lower bound give $\neg(\mathrm{SAT}\in\mathrm{P})$, i.e.\ $\mathrm{P}\ne\mathrm{NP}$ by Cook--Levin.
($\Leftarrow$) If $\mathrm{P}\ne\mathrm{NP}$ the antecedent is false, so the implication holds.
(b) PV $\vdash\forall F\,[A(F)=\mathrm{UNSAT}\to\neg\mathrm{Sat}(F,\cdot)]$ translates (Cook 1975) into polynomial-size EF proofs of the propositional
soundness statements $\llbracket\cdot\rrbracket_N$. Composing with an EF evaluation of the run of $A$ on the fixed input $F$ (polynomial size, since EF extension
variables can name every gate of the run) refutes $F$ in polynomial size.
\end{proof}

\begin{corollary}
The ``weakest sufficient bridge'' is never strictly weaker than the full problem for any calculus whose lower bound we can currently prove.
For the strong calculus EF the bridge is available only for PV-verified algorithms, and closing the target there would also need a superpolynomial EF lower bound,
a long-standing open problem. The combination ``EF superpolynomial $\Rightarrow$ no PV-verified polynomial SAT algorithm'' is a known consequence of Cook 1975;
Pich--Santhanam (arXiv:2312.08163) give additional hypotheses under which EF hardness implies $\mathrm{P}\ne\mathrm{NP}$ outright.
\end{corollary}

\section{What free choice of the optimal sequence buys}
Granting the best reduction sequence for free is legitimate for a lower bound: $\min_\pi\operatorname{cost}(\pi)\le\operatorname{cost}(\pi_A)$
for any particular derivation $\pi_A$ extracted from an algorithm. The consequence is that $\WR{R}$ measures \emph{nondeterministic} certificate size:
a lower bound on it is a lower bound against coNP $\subseteq$ NP via $R$, not directly against deterministic time. The deterministic
content of P vs NP enters only through the bridge, which must convert the trace of an arbitrary deterministic algorithm into an $R$-derivation.
A computation step of an arbitrary algorithm is a Boolean circuit gate; a calculus that can absorb such steps at polynomial overhead must
be able to define new variables for arbitrary gates, which is exactly the extension rule. So the calculus either lacks extension
(horn (a): lower bounds provable, bridge $\equiv$ P$\ne$NP) or has it (horn (b): bridge conditional, lower bound open).

\section{Multisign DAG hypergraph calculus}\label{s:hyper}
Encode a 3-CNF as a signed hypergraph (vertices $=$ variables, hyperedges $=$ clauses with sign vectors) and let a \emph{reduction step}
be any sound local rule merging or resolving hyperedges, deleting satisfied structure, or recording a shared DAG node. Then:
\begin{itemize}
\item If every rule has bounded arity and each new hyperedge is a clause (or a bounded-width conjunction of clauses), the calculus is p-simulated
by resolution or Res$(k)$ and the minimal reduced state is the empty hyperedge. Then $\WR{R}$ is superpolynomial on Tseitin formulas of expanders,
random 3-CNF with $\Delta n$ clauses, and the pigeonhole principle: KNOWN THEOREM by transfer. By Thm.~\ref{t:dilemma}(a), the bridge is then P$\ne$NP itself.
\item If the calculus may introduce named DAG nodes for arbitrary Boolean functions of existing nodes (shared DAGs, circuit compression), it is equivalent to EF (or
to Frege if names are formulas). No lower bound is known: OPEN.
\item Adding GF(2) linear rows (parity reasoning) defeats Tseitin hardness: Gaussian elimination refutes Tseitin in polynomial size, and
Res(lin)/XOR-resolution lower bounds are known only in restricted (e.g.\ tree-like or regular) forms. Hard families surviving algebraic reasoning are then random 3-CNF and
PHP-type families (polynomial calculus lower bounds), but none is known to survive EF.
\item Adaptive strategy switching does not escape horn (a): switching among sub-calculi is again a single calculus (their union), and the union's
lower bound is the minimum over its parts' simulations.
\end{itemize}
Status: CONDITIONAL on the precise rule list; the repository's R\_mix341 authentication (proof2.txt P2-D1) decides which horn it lies in.
Note: R\_mix341 includes GF(2) and paid auxiliary definitions; if those definitions may name arbitrary gates, R\_mix341 is EF-like (horn (b)).

\section{Provable algorithm-class bridges}
\begin{proposition}[KNOWN THEOREM]\label{p:cdcl}
Every run of CDCL with restarts and arbitrary decision/learning heuristics that refutes $F$ in $t$ steps yields a resolution refutation of size $\mathrm{poly}(t)$
(Beame--Kautz--Sabharwal 2004). Hence CDCL needs superpolynomial time on Tseitin, PHP and random 3-CNF.
\end{proposition}
The productive version of the objective is therefore: \textbf{for a named algorithm class $\mathcal A$ prove Bridge$_{\mathcal A}(R)$, then use the known $\WR{R}$ lower bound.}
Classes with this property: DPLL (tree-like resolution), CDCL (resolution), Gröbner-based solvers (polynomial calculus), ILP with Chvátal cuts (cutting planes).
Each gives an unconditional lower bound for that class only. Extending $\mathcal A$ to ``all polynomial-time algorithms'' is horn (b).

\section{Barrier-transfer and literature preflight}
Main obstacles: (O1) no superpolynomial EF/Frege lower bound; (O2) no unconditional simulation of arbitrary algorithms by a fixed calculus.
Literature checked this run: Cook--Reckhow 1979; Cook 1975; Krajíček--Pudlák 1989/1990; Pich--Santhanam, \emph{Towards P$\ne$NP from Extended Frege lower bounds}, arXiv:2312.08163 (2023);
Cook--Krajíček, JSL 2007 (consistency of P$\ne$NP with bounded arithmetic from EF lower bounds); the 2025 automatability preprint arXiv:2506.16956 (proof analysis for EF is NP-complete).
Transfer domains: control theory \emph{barrier certificates} (existence of a certificate vs.\ synthesizing it) --- ANALOGY: same gap as
proof existence vs.\ algorithmic search, no theorem transfers. Coding theory / PQ items 1--3 (Schur squares, Hilbert functions) --- FAILED TRANSFER for this question:
they are algebraic invariants of instances, while the bridge is about simulation of algorithms. PQ item 5 (LWE-style worst/average reduction) --- ANALOGY: an average-case
$\WR{R}$ lower bound (random 3-CNF) already exists for weak $R$; no worst-case-to-algorithm reduction known.
Relativization: Prop.~\ref{p:RA} and Thm.~\ref{t:dilemma}(a) relativize; they create no new barrier, but they show the bridge inherits every barrier of P vs NP.

\section{Corpus reviewed}
Root AGENTS.md; proof1.txt and proof2.txt (heads, P2-D1 R\_mix341 entry); root main.tex (size and status only, 1.9\,MB);
latest LSPV-1..5 activity on 2026-11-10 (git log: LSPV-1/556, LSPV-2/541, LSPV-3/553--554, LSPV-4/526);
coordination handoffs mentioning worst-case quotients (LSPV1 MUS\_WORSTCASE\_QUOTIENT350); Claude world ledger rows L01--L40.

\section{Duplication / contradiction audit}
Prior LSPV reports already prove restricted lower bounds for R\_mix341 fragments and model-boundary controls (proof2.txt P2-D1 evidence, LSPV-3/464:
the external FRAT-XOR grammar refutes Tseitin in $h+2$ records). That agrees with \S\ref{s:hyper} item 3. No prior report found in this survey stated the
bridge dilemma or the equivalence Bridge$(R)\iff$P$\ne$NP for lower-bounded $R$; this is a standard observation in proof complexity, not a new theorem.
No contradiction with root main.tex was found; main.tex already keeps P vs NP OPEN.

\section{What changed relative to root main}
Nothing in root main.tex is edited (this lane owns no region). New: world ledger rows L41--L45 and this report.

\section{Adopted / rejected / retired components}
Adopted: free optimal-sequence minimization (sound for lower bounds). Rejected as a path to P$\ne$NP: any bridge for a calculus without unrestricted extension,
since it is equivalent to P$\ne$NP. Not retired: horn (b) via EF and the Pich--Santhanam hypotheses.

\section{Stale or overstated claims found}
The shared directive calls Bridge$(R)$ the ``weakest sufficient bridge''. For weak $R$ it is not weaker than P$\ne$NP (Thm.~\ref{t:dilemma}(a)). Recommend
recording it as a reformulation (world ledger L42: ``Equivalent to P$\ne$NP'').

\section{Next targets}
\begin{enumerate}
\item Decide which horn R\_mix341 lies in: can its paid auxiliary definitions name arbitrary gates? If yes, $\WR{R}$ lower bounds for it are EF-hard.
\item Prove Bridge$_{\mathcal A}$(R) for a broader explicit algorithm class (e.g.\ CDCL $+$ Gaussian elimination $\to$ Res(lin)) and pair with the best known Res(lin) lower bounds.
\item Survey whether any hard family is a candidate for EF hardness (e.g.\ random 3-CNF, Razborov's pseudorandom-generator tautologies) and record it as OPEN.
\end{enumerate}

\noindent\textbf{P versus NP remains OPEN.} Every load-bearing arrow of the worst-case spine (a superpolynomial $\WR{R}$ for an EF-strength $R$, and Bridge$(R)$ for all polynomial-time algorithms) is OPEN.

\end{document}
% END_OF_VERIFIED_REPORT::CLAUDE::WORSTCASE-BRIDGE-0001
